\section{Exercício prático: WordPress + MySQL com Docker Compose}

Neste exercício, o objetivo é montar um ambiente completo de WordPress (aplicação + banco MySQL) usando Docker Compose, configurando a rede para que eles se comuniquem automaticamente.

\subsection{Estrutura do projeto}

\begin{verbatim}
wordpress-docker/
  compose.yaml
\end{verbatim}

\subsection{Arquivo compose.yaml}

\begin{verbatim}
services:
  db:
    image: mysql:8.0
    restart: always
    environment:
      MYSQL_ROOT_PASSWORD: senharoot
      MYSQL_DATABASE: wordpress
      MYSQL_USER: wpuser
      MYSQL_PASSWORD: wpsenha
    volumes:
      - db_data:/var/lib/mysql
    networks:
      - wp-network

  wordpress:
    image: wordpress:latest
    restart: always
    ports:
      - "8080:80"
    environment:
      WORDPRESS_DB_HOST: db:3306
      WORDPRESS_DB_USER: wpuser
      WORDPRESS_DB_PASSWORD: wpsenha
      WORDPRESS_DB_NAME: wordpress
    volumes:
      - wp_data:/var/www/html
    networks:
      - wp-network
    depends_on:
      - db

volumes:
  db_data:
  wp_data:

networks:
  wp-network:
\end{verbatim}

\subsection{Explicação do arquivo}

\subsubsection{Serviço db (MySQL)}

\begin{itemize}
    \item Usa a imagem \texttt{mysql:8.0}.
    \item As variáveis de ambiente configuram o banco: senha root, nome do banco, usuário e senha.
    \item O volume \texttt{db\_data} persiste os dados do MySQL em \texttt{/var/lib/mysql}, garantindo que os dados sobrevivam a reinicializações.
    \item Conectado à rede \texttt{wp-network}.
\end{itemize}

\subsubsection{Serviço wordpress}

\begin{itemize}
    \item Usa a imagem oficial \texttt{wordpress:latest}.
    \item Porta 8080 do host mapeada para porta 80 do container.
    \item As variáveis de ambiente apontam para o serviço \texttt{db} --- o Docker Compose resolve o nome \texttt{db} automaticamente para o IP do container MySQL, graças à resolução DNS da rede customizada.
    \item O volume \texttt{wp\_data} persiste os arquivos do WordPress.
    \item \texttt{depends\_on: db} garante que o MySQL inicie antes do WordPress.
    \item Conectado à mesma rede \texttt{wp-network}.
\end{itemize}

\subsubsection{Volumes e Networks}

Os volumes \texttt{db\_data} e \texttt{wp\_data} são declarados na seção \texttt{volumes} para que o Docker os gerencie. A rede \texttt{wp-network} é declarada na seção \texttt{networks} e permite que os dois serviços se comuniquem pelo nome do serviço (por exemplo, o WordPress acessa o MySQL usando o hostname \texttt{db}).

\subsection{Executando}

\begin{verbatim}
# Subir toda a stack
docker compose up -d

# Verificar os serviços rodando
docker compose ps

# Acessar o WordPress no navegador
# http://localhost:8080
\end{verbatim}

Ao acessar \texttt{http://localhost:8080}, o assistente de instalação do WordPress será exibido. Todo o ambiente --- aplicação e banco de dados --- foi iniciado com um único comando.

\subsection{Verificando a comunicação entre containers}

\begin{verbatim}
# Entrar no container do WordPress
docker exec -it wordpress-docker-wordpress-1 bash

# Dentro do container, testar a conexão com o MySQL
apt-get update && apt-get install -y default-mysql-client
mysql -h db -u wpuser -pwpsenha wordpress -e "SHOW TABLES;"
\end{verbatim}

O hostname \texttt{db} é resolvido automaticamente pelo Docker DNS para o IP do container MySQL, sem necessidade de configuração manual.

\subsection{Derrubando o ambiente}

\begin{verbatim}
# Parar e remover containers (volumes são preservados)
docker compose down

# Parar, remover containers E volumes (apaga tudo)
docker compose down -v
\end{verbatim}

A flag \texttt{-v} no \texttt{docker compose down} remove também os volumes, apagando todos os dados. Sem ela, os volumes são preservados e os dados estarão lá na próxima vez que subir o ambiente.
